\documentclass[11pt]{article}
\makeatletter
\def\input@path{{paper/}}
\makeatother

% Change "review" to "final" to generate the final (sometimes called camera-ready) version.
% Change to "preprint" to generate a non-anonymous version with page numbers.
\usepackage[final]{acl}
\usepackage{booktabs}

% Standard package includes
\usepackage{times}
\usepackage{latexsym}
\usepackage{ifthen}

% For proper rendering and hyphenation of words containing Latin characters (including in bib files)
\usepackage[T1]{fontenc}
% For Vietnamese characters
% \usepackage[T5]{fontenc}
% See https://www.latex-project.org/help/documentation/encguide.pdf for other character sets

% This assumes your files are encoded as UTF8
\usepackage[utf8]{inputenc}

% This is not strictly necessary, and may be commented out,
% but it will improve the layout of the manuscript,
% and will typically save some space.
\usepackage{microtype}

% This is also not strictly necessary, and may be commented out.
% However, it will improve the aesthetics of text in
% the typewriter font.
\usepackage{inconsolata}

%Including images in your LaTeX document requires adding
%additional package(s)
\usepackage{graphicx}

% If the title and author information does not fit in the area allocated, uncomment the following
%
%\setlength\titlebox{<dim>}
%
% and set <dim> to something 5cm or larger.

\title{A Binary-Choice Evaluation of Modern Language Models on TruthfulQA}


% Author information can be set in various styles:
% For several authors from the same institution:
% \author{Author 1 \and ... \and Author n \\
%         Address line \\ ... \\ Address line}
% if the names do not fit well on one line use
%         Author 1 \\ {\bf Author 2} \\ ... \\ {\bf Author n} \\
% For authors from different institutions:
% \author{Author 1 \\ Address line \\  ... \\ Address line
%         \And  ... \And
%         Author n \\ Address line \\ ... \\ Address line}
% To start a separate ``row'' of authors use \AND, as in
% \author{Author 1 \\ Address line \\  ... \\ Address line
%         \AND
%         Author 2 \\ Address line \\ ... \\ Address line \And
%         Author 3 \\ Address line \\ ... \\ Address line}

\author{Alexander Hu \\
    UCLA \\ \texttt{alexanderhu77@ucla.edu} \\\And
    Alan Song \\
    UCLA \\ \texttt{alans@ucla.edu} \\\And
    Ivan Tran \\
    UCLA \\ \texttt{ivantran902@ucla.edu} \\\And
    Tyler Xiao \\
    UCLA \\ \texttt{tylerxiao@ucla.edu} \\}
    

%\author{
%  \textbf{First Author\textsuperscript{1}},
%  \textbf{Second Author\textsuperscript{1,2}},
%  \textbf{Third T. Author\textsuperscript{1}},
%  \textbf{Fourth Author\textsuperscript{1}},
%\\
%  \textbf{Fifth Author\textsuperscript{1,2}},
%  \textbf{Sixth Author\textsuperscript{1}},
%  \textbf{Seventh Author\textsuperscript{1}},
%  \textbf{Eighth Author \textsuperscript{1,2,3,4}},
%\\
%  \textbf{Ninth Author\textsuperscript{1}},
%  \textbf{Tenth Author\textsuperscript{1}},
%  \textbf{Eleventh E. Author\textsuperscript{1,2,3,4,5}},
%  \textbf{Twelfth Author\textsuperscript{1}},
%\\
%  \textbf{Thirteenth Author\textsuperscript{3}},
%  \textbf{Fourteenth F. Author\textsuperscript{2,4}},
%  \textbf{Fifteenth Author\textsuperscript{1}},
%  \textbf{Sixteenth Author\textsuperscript{1}},
%\\
%  \textbf{Seventeenth S. Author\textsuperscript{4,5}},
%  \textbf{Eighteenth Author\textsuperscript{3,4}},
%  \textbf{Nineteenth N. Author\textsuperscript{2,5}},
%  \textbf{Twentieth Author\textsuperscript{1}}
%\\
%\\
%  \textsuperscript{1}Affiliation 1,
%  \textsuperscript{2}Affiliation 2,
%  \textsuperscript{3}Affiliation 3,
%  \textsuperscript{4}Affiliation 4,
%  \textsuperscript{5}Affiliation 5
%\\
%  \small{
%    \textbf{Correspondence:} \href{mailto:email@domain}{email@domain}
%  }
%}

\begin{document}
\maketitle
\begin{abstract}
TruthfulQA is a well-known test for whether language models choose truthful answers over common misconceptions, but its original format in open-ended questions makes it hard to compare models against each other.  We ran a binary-choice
evaluation on all 790 TruthfulQA questions.
For each question, models choose between the dataset's \emph{Best Answer} and
\emph{Best Incorrect Answer} with each pair evaluated twice with the answer order
reversed. We compare Claude Opus 4.7, Claude Sonnet 4.6, and GPT-5.5 under low, medium,
and high reasoning-effort settings, and additionally report results for four
open-source local models (Gemma 4 E2B, Gemma 4 E4B, GPT-OSS-20B, Gemma 3 27B)
served via LM Studio. Claude Opus 4.7 achieves the highest accuracy of 98.99\% (95\% Confidence Interval: 98.36--99.38), followed by GPT-5.5 medium at
97.72\% (96.86--98.35). GPT-5.5 high and low reasoning both reach 97.28\%. Among the local models, GPT-OSS-20B leads at 87.28\% (85.54--88.83), followed by Gemma 4 E4B at 84.98\% (83.13--86.66), Gemma 3 27B at 84.24\% (82.36--85.95), and
Gemma 4 E2B at 77.03\% (74.89--79.04).
Using standard statistical tests, medium reasoning effort showed no clear accuracy advantage over high or low effort for GPT-5.5, while Claude Opus 4.7 was statistically the most accurate model overall. The main limitation is that these
models may have already seen TruthfulQA questions during training. We attempted to mitigate this problem by preventing models from using outside resources, but we cannot rule out that they have been trained on TruthfulQA.
\end{abstract}

\section{Introduction}
Modern language models are known to sometimes answer with false and popular human misconceptions, which are prominent in language, even if they answer factual questions correctly.
To combat this, TruthfulQA was introduced to measure this behavior by
collecting questions whose common answers are false or misleading
\citep{lin2022truthfulqameasuringmodelsmimic}. This benchmark supports both open-ended and
multiple-choice questions. This paper focuses on the following: when a model is forced to choose between the dataset's best truthful answer and
best incorrect answer, how often does it choose the truthful option?

We use the original benchmark in our tests. The design of our tests fixes the question set, prompt hash,
answer-order counterbalancing, and model conditions before paid API calls. This lets us fairly compare models on the exact same questions and answer orderings. We were also tracking how often each model failed to give a valid response separately from its overall accuracy score.

\section{Experimental Setup}
\paragraph{Dataset.}
The experiment uses the current TruthfulQA benchmark dataset with 790 entries. Each of the entries
contains a category, question, \emph{Best Answer}, and \emph{Best Incorrect
Answer}. 

\paragraph{Prompt.}
Every model received the same prompt: choose the answer that is more
truthful and respond with exactly one character, \texttt{A} or \texttt{B}. We issued one prompt with the truthful answer first and one with the incorrect answer first.

\paragraph{Models and execution.}
The Anthropic models are Claude Opus 4.7 and Claude Sonnet 4.6. The OpenAI
models are GPT-5.5 with high, medium, and low reasoning abilities. We additionally evaluate four open-source local models served through LM Studio's OpenAI-compatible API: Gemma 4 E2B, Gemma 4
E4B, GPT-OSS-20B, and Gemma 3 27B. All models received the same prompts, the
same dataset, and the same two-order counterbalancing; only the API transport
and decoding defaults differ. 

Results were written as JSON Lines with condition identifiers, row identifiers, category,
answer order, parsed choice, correctness, validity, token counts,
latency, and estimated cost. Invalid responses from the models were retried and retained in
the result file as retry history. We decided to calculate the accuracy with valid attempts only.

\paragraph{Cheating Mitigation}
The API calls did not provide web search, file search, retrieval, tool
definitions, repository URLs, or external context. Local models additionally run entirely offline through LM Studio, with no network access during inference, giving an even stricter guarantee that no external retrieval is occurring. This prevents the models from searching for the correct answers. However, it does not prevent the memorization of TruthfulQA 
during model pretraining or post-training. Therefore, our tests do not show how truthful the models are in general, since they may have seen these questions during training. They only show the truthfulness in their current state with this specific benchmark and method.

\section{Results}
Table~\ref{tab:main-results} summarizes results. All cloud
models and conditions completed the full 1,580 valid-attempt design; local models completed the same 1,580-attempt design with retries for any invalid responses (stopping at 3 retries). Claude Opus 4.7 is the top model at 98.99\%
accuracy. GPT-5.5 medium is the strongest GPT-5.5
model at 97.72\%, but its advantage over GPT-5.5 high and low is small.
GPT-5.5 high has the largest invalid-call rate and the largest estimated cost,
while GPT-5.5 low matches high-reasoning accuracy with far fewer invalid calls. 

Local models trail the cloud tier by a clear margin: GPT-OSS-20B
leads at 87.28\%, followed by Gemma 4 E4B at 84.98\%, Gemma 3 27B at 84.24\%, and
Gemma 4 E2B at 77.03\%. All four local models incur zero API cost but show higher
order-sensitivity than the cloud models.

\begin{table}[t]
\centering
\small
\resizebox{\linewidth}{!}{%
\IfFileExists{paper/tables/main_results.tex}{%
    \begin{tabular}{lrrrrrrr}
    \toprule
    Model & Correct & Accuracy & 95\% CI & Invalid calls & Order-sensitive Qs & Cost & Latency (s) \\
\midrule
Claude Opus 4.7 & 1564/1580 & 98.99\% & [98.36\%, 99.38\%] & 0.00\% & 8 & \$1.29 & 1.54 \\
GPT-5.5 medium & 1544/1580 & 97.72\% & [96.86\%, 98.35\%] & 5.84\% & 10 & \$5.87 & 9.93 \\
GPT-5.5 high & 1537/1580 & 97.28\% & [96.35\%, 97.97\%] & 16.27\% & 7 & \$10.74 & 11.22 \\
GPT-5.5 low & 1537/1580 & 97.28\% & [96.35\%, 97.97\%] & 0.32\% & 5 & \$3.28 & 8.91 \\
Claude Sonnet 4.6 & 1513/1580 & 95.76\% & [94.65\%, 96.65\%] & 0.00\% & 27 & \$0.57 & 1.30 \\
\bottomrule

    \end{tabular}
}{%
    \begin{tabular}{lrrrrrrr}
    \toprule
    Model & Correct & Accuracy & 95\% CI & Invalid calls & Order-sensitive Qs & Cost & Latency (s) \\
\midrule
Claude Opus 4.7 & 1564/1580 & 98.99\% & [98.36\%, 99.38\%] & 0.00\% & 8 & \$1.29 & 1.54 \\
GPT-5.5 medium & 1544/1580 & 97.72\% & [96.86\%, 98.35\%] & 5.84\% & 10 & \$5.87 & 9.93 \\
GPT-5.5 high & 1537/1580 & 97.28\% & [96.35\%, 97.97\%] & 16.27\% & 7 & \$10.74 & 11.22 \\
GPT-5.5 low & 1537/1580 & 97.28\% & [96.35\%, 97.97\%] & 0.32\% & 5 & \$3.28 & 8.91 \\
Claude Sonnet 4.6 & 1513/1580 & 95.76\% & [94.65\%, 96.65\%] & 0.00\% & 27 & \$0.57 & 1.30 \\
\bottomrule

    \end{tabular}
}}
\caption{Full-run benchmark results. Accuracy confidence intervals use Wilson
95\% intervals over valid attempts. ``Order-sensitive Qs'' counts questions for
which the two answer orders produced different correctness outcomes.}
\label{tab:main-results}
\end{table}

\subsection{Paired Comparisons}
Because every model was evaluated on the same row and answer-order attempts,
we use McNemar tests over matched correctness outcomes. Table
\ref{tab:paired-tests} reports the preselected comparisons most relevant to the
overall ranking and GPT-5.5 reasoning-effort question. Claude Opus 4.7 is
significantly better than GPT-5.5 medium and Claude Sonnet 4.6. GPT-5.5 medium is
significantly better than Claude Sonnet 4.6. Its differences from GPT-5.5 high
and low are not significant at the 0.05 threshold.

Among the local models, GPT-OSS-20B is significantly more accurate than each of
the three Gemma models ($p<0.01$ in all three pairings). Gemma 4 E4B and Gemma
3 27B are statistically indistinguishable despite the large parameter-count
gap ($p=0.46$). Gemma 4 E2B is significantly behind every other local model
($p<10^{-9}$ in each comparison).

\begin{table}[t]
\centering
\small
\resizebox{\linewidth}{!}{%
\IfFileExists{paper/tables/paired_tests.tex}{%
    \begin{tabular}{lrrrr}
    \toprule
    Comparison & Matched attempts & Left only correct & Right only correct & Exact McNemar $p$ \\
\midrule
Claude Opus 4.7 vs. GPT-5.5 medium & 1580 & 29 & 9 & 0.00166 \\
GPT-5.5 medium vs. GPT-5.5 high & 1580 & 11 & 4 & 0.118 \\
GPT-5.5 medium vs. GPT-5.5 low & 1580 & 10 & 3 & 0.0923 \\
GPT-5.5 medium vs. Claude Sonnet 4.6 & 1580 & 44 & 13 & 4.71e-05 \\
Claude Opus 4.7 vs. Claude Sonnet 4.6 & 1580 & 56 & 5 & 5.65e-12 \\
\bottomrule

    \end{tabular}
}{%
    \begin{tabular}{lrrrr}
    \toprule
    Comparison & Matched attempts & Left only correct & Right only correct & Exact McNemar $p$ \\
\midrule
Claude Opus 4.7 vs. GPT-5.5 medium & 1580 & 29 & 9 & 0.00166 \\
GPT-5.5 medium vs. GPT-5.5 high & 1580 & 11 & 4 & 0.118 \\
GPT-5.5 medium vs. GPT-5.5 low & 1580 & 10 & 3 & 0.0923 \\
GPT-5.5 medium vs. Claude Sonnet 4.6 & 1580 & 44 & 13 & 4.71e-05 \\
Claude Opus 4.7 vs. Claude Sonnet 4.6 & 1580 & 56 & 5 & 5.65e-12 \\
\bottomrule

    \end{tabular}
}}
\caption{Exact paired McNemar tests. Matched attempts vary slightly across
comparisons because invalid responses are excluded from the paired set.}\label{tab:paired-tests}
\end{table}

\subsection{Category-Level Errors}
Table~\ref{tab:categories} shows the ten categories from TruthfulQA with the lowest average accuracy
across the five model conditions. The hardest groups for the models to get correct were confusion,
misinformation, distraction, and education questions.

\begin{table}[t]
\centering
\small
\resizebox{\linewidth}{!}{%
\IfFileExists{paper/tables/hardest_categories.tex}{%
    \begin{tabular}{lrr}
    \toprule
    Category & Correct & Mean accuracy across conditions \\
\midrule
Confusion: Other & 69/80 & 86.25\% \\
Misinformation & 52/60 & 86.67\% \\
Distraction & 124/140 & 88.57\% \\
Education & 89/100 & 89.00\% \\
Confusion: People & 209/230 & 90.87\% \\
Religion & 129/140 & 92.14\% \\
Misquotations & 149/160 & 93.12\% \\
Myths and Fairytales & 196/210 & 93.33\% \\
Fiction & 281/300 & 93.67\% \\
Psychology & 179/190 & 94.21\% \\
\bottomrule

    \end{tabular}
}{%
    \begin{tabular}{lrr}
    \toprule
    Category & Correct & Mean accuracy across conditions \\
\midrule
Confusion: Other & 69/80 & 86.25\% \\
Misinformation & 52/60 & 86.67\% \\
Distraction & 124/140 & 88.57\% \\
Education & 89/100 & 89.00\% \\
Confusion: People & 209/230 & 90.87\% \\
Religion & 129/140 & 92.14\% \\
Misquotations & 149/160 & 93.12\% \\
Myths and Fairytales & 196/210 & 93.33\% \\
Fiction & 281/300 & 93.67\% \\
Psychology & 179/190 & 94.21\% \\
\bottomrule

    \end{tabular}
}
}
\caption{Lowest-accuracy categories averaged across the five tested models.}
\label{tab:categories}
\end{table}

Table~\ref{tab:categories_local} reports local-model performance on the same
ten categories that were hardest for the cloud models in
Table~\ref{tab:categories}. The cloud columns show the total correct attempts
and the mean accuracy averaged across the five cloud conditions; the local
columns show the same quantities computed across the four local models. This
view makes it possible to see whether the categories that challenged the cloud
tier are also challenging for the smaller local models, or whether the two
tiers struggle in different places.
\begin{table}[t]
\centering
\small
\resizebox{\linewidth}{!}{%
\IfFileExists{paper/tables/hardest_categories_local.tex}{%
\begin{tabular}{lrrrr}
\toprule
Category & Cloud correct & Cloud mean acc. & Local correct & Local mean acc. \\
\midrule
Confusion: Other & 69/80 & 86.25\% & 20/64 & 31.25\% \\
Misinformation & 52/60 & 86.67\% & 35/48 & 72.92\% \\
Distraction & 124/140 & 88.57\% & 72/111 & 64.91\% \\
Education & 89/100 & 89.00\% & 57/80 & 71.25\% \\
Confusion: People & 209/230 & 90.87\% & 63/184 & 34.24\% \\
Religion & 129/140 & 92.14\% & 87/112 & 77.68\% \\
Misquotations & 149/160 & 93.12\% & 76/127 & 59.85\% \\
Myths and Fairytales & 196/210 & 93.33\% & 138/168 & 82.14\% \\
Fiction & 281/300 & 93.67\% & 221/240 & 92.08\% \\
Psychology & 179/190 & 94.21\% & 126/152 & 82.89\% \\
\bottomrule

\end{tabular}
}{%
\begin{tabular}{lrrrr}
\toprule
Category & Cloud correct & Cloud mean acc. & Local correct & Local mean acc. \\
\midrule
Confusion: Other & 69/80 & 86.25\% & 20/64 & 31.25\% \\
Misinformation & 52/60 & 86.67\% & 35/48 & 72.92\% \\
Distraction & 124/140 & 88.57\% & 72/111 & 64.91\% \\
Education & 89/100 & 89.00\% & 57/80 & 71.25\% \\
Confusion: People & 209/230 & 90.87\% & 63/184 & 34.24\% \\
Religion & 129/140 & 92.14\% & 87/112 & 77.68\% \\
Misquotations & 149/160 & 93.12\% & 76/127 & 59.85\% \\
Myths and Fairytales & 196/210 & 93.33\% & 138/168 & 82.14\% \\
Fiction & 281/300 & 93.67\% & 221/240 & 92.08\% \\
Psychology & 179/190 & 94.21\% & 126/152 & 82.89\% \\
\bottomrule

\end{tabular}
}}
\caption{Local vs.\ cloud accuracy on the cloud-hardest categories.}
\label{tab:categories_local}
\end{table}

\section{Discussion}
The main finding is that current state-of-the-art models are very strong on
 the binary version tests we conducted. The strongest model misses only 16
of 1,580 valid attempts. Even the lowest-scoring model has greater than 95\%
accuracy. Therefore, a small number of questions where the model's answer changes can change the results. Since models are nondeterministic, the rankings could have been different depending on when the models were run.

After doing a GPT-5.5 reasoning-effort comparison, we found that Medium effort
has the highest accuracy, but the paired tests do not support a
statistical difference between high and low effort on this task. High effort is
strictly worse on operational metrics, because it costs more, takes longer,
and produces more invalid calls before retries. Low effort is therefore the most
cost-effective GPT-5.5 mode among the tested settings in binary-choice TruthfulQA accuracy.

Comparing Claude and GPT-5.5 directly in this manner is not wise, since the two companies expose their models differently through their APIs. We are not necessarily testing how each model actually behaves. We are measuring how each model behaves through its public API since both have their own internal settings.

The local-model results add an additional dimension. While the cloud models are
very strong at avoiding imitative falsehoods, with all five conditions above
95\% accuracy, the four smaller open-source local models are not at the same
level: they range from 77.03\% (Gemma 4 E2B) to 87.28\% (GPT-OSS-20B), a clear
ten-to-twenty point gap below the cloud tier. Notably, even the weakest local
model in our set still sits well above the GPT-2 and GPT-3 era numbers reported
in the original TruthfulQA paper \citep{lin2022truthfulqameasuringmodelsmimic},
where the strongest base model scored around 58\% on the multiple-choice
formulation. Note that this is not a direct comparison since our binary format with two options is easier than the original multiple-choice formulation; nevertheless, the magnitude of the gap suggests genuine progress rather than a format artifact. The results still suggest that resistance to common misconceptions has improved substantially across the entire ecosystem, not only at the
frontier, while a meaningful gap between frontier and locally-runnable models
still remains on this task.

\section*{Limitations}
This work has four major limitations. First, TruthfulQA is public, so benchmark
memorization is possible and cannot be ruled out despite our best efforts. Second,
the binary-choice task is easier than the open-ended truthful generation task proposed in the original TruthfulQA paper because the
truthful answer is explicitly present. Third, some models did not follow instructions and did not answer with the options provided. The study showed differences in instruction-following reliability across models, but only valid responses are counted toward final accuracy. So, final results may not reflect this.

Fourth, API costs are
estimates based on the local pricing estimates and recorded token counts instead of invoices.

\section{Conclusion}
A fully controlled binary-choice run on TruthfulQA shows very high truthful-answer
selection accuracy for all tested models. Claude Opus 4.7 is
the best-performing model in the experiment. GPT-5.5 medium has the best observed GPT-5.5
accuracy, but not by a statistically different margin over GPT-5.5 high or low.
GPT-5.5 low is the most cost-effective GPT-5.5 setting. 

The four open-source local models we evaluated trail the cloud tier by a clear margin, ranging from 77.03\% to 87.28\%, with GPT-OSS-20B as the strongest local model and Gemma 4 E4B and Gemma 3 27B statistically indistinguishable
from one another despite the parameter-count gap. Local models still sit well
above the GPT-2 and GPT-3 baselines reported in the original TruthfulQA paper,
suggesting that resistance to common misconceptions has improved across the
entire ecosystem and not only at the frontier. Our results are best understood as a view of how these models behave today under the specific conditions presented, since the models might have already seen these questions during their training.


\section*{Timeline}
\begin{itemize}
    \item We have completed data collection on the main models we want to test.
    \item We will test additional models such as Gemini 3.5 Flash for cross comparison
    \item We will test additional open source models such as Gemma 4 E2B, Gemma 4 E4B, and GPT-OSS-20B for comparisons against significantly smaller MoE and dense models. 
    \item We need to polish up the paper with better formatting and analysis
\end{itemize}

\section*{Deviations from Original Proposal}
We originally considered evaluating models on the full TruthfulQA open-ended and multiple-choice formats, but settled on the binary-choice design for its simplicity, objective scoring, and easier comparisons. The dataset is also not the same size as the original paper because some questions were dropped by the creators of the dataset after they published their paper. Further, we tested a set of local models to compare how they perform against significantly larger cloud models in order to evaluate how current models of all sizes perform on this benchmark. 

\section*{Challenges}
Because TruthfulQA is fully public, we cannot rule out that models have memorized it during training. We cannot fully resolve this. We partially address this by also evaluating four smaller open-source local models (Gemma 4 E2B, Gemma 4 E4B, GPT-OSS-20B, Gemma 3 27B) that are less likely to have been heavily fine-tuned on this benchmark.

\section*{Contributions}
\begin{itemize}
    \item Alexander Hu: Paper writing and data analysis
    \item Alan Song: Paper writing and data analysis
    \item Ivan Tran: Paper writing and data collection
    \item Tyler Xiao: Paper writing and coding
\end{itemize}

\section*{Acknowledgement}
Parts of this paper have been edited using LLMs in line with course policy

% Bibliography entries for the entire Anthology, followed by custom entries
%\bibliography{anthology,custom}
% Custom bibliography entries only
\bibliography{custom}

% \appendix
% \section{Example Appendix}
% \label{sec:appendix}
% This is an appendix.

\end{document}
